\chapter{Background}\label{ch:background}

\section{Lexical linear models}
In \textit{Word TF-IDF LR}, each article is represented using a fitted vocabulary of up to 50,000 single words (unigrams) and adjacent word pairs (bigrams). The model can use a single word such as \texttt{market} on its own and give extra weight to an informative phrase such as \texttt{stock market}. As \citet{wang2012} show, pairing single words with local word pairs forms a strong baseline for topic and sentiment tasks. Because a short AG News story uses only a few dozen distinct terms, almost every entry in its 50,000-slot vector stays zero, making the model lightweight and fast.

To stop long summaries or frequent filler words from overpowering concise headlines, weights follow TF-IDF scaling and unit length normalization \citep{tfidf}. Logarithmic term-frequency scaling reduces the impact of repeating the same word multiple times in one story, inverse document frequency (IDF) shrinks weights for generic words shared across topics, and Euclidean normalization scales every article vector to length one so proportions matter more than story length. Multinomial logistic regression then computes a weighted sum for each topic with a smooth $L_2$ weight penalty that spreads credit across related clues \citep{logisticregression}.

When a spelling swap hits one word, that word and its two adjacent word pairs drop out of the checklist. If the same word appears cleanly elsewhere in the story, it still votes, and dozens of intact words across the headline and summary keep pointing toward the right category.

\section{Neural text models}
Following the text CNN design of \citet{kim2014}, \textit{Word CNN} preserves word order by mapping each known vocabulary word to a learned 96-dimensional vector initialized randomly. Words absent from the training dictionary, including nearly any newly created typo, map to one shared unknown-word vector.

Three parallel filter banks slide across three, four, and five consecutive words at a time, looking for short topical phrases anywhere in the article. Each filter passes through a rectified linear unit (ReLU), and global max pooling keeps only each filter's single highest score across valid article windows before a linear layer scores the four topics.

Keeping only the single strongest match per filter works well when one sharp phrase defines the story, yet makes the prediction sensitive to where a typo lands. If a letter swap breaks a word inside that peak phrase, the filter immediately drops to a weaker backup match elsewhere in the text, even when neighboring words remain readable. Comparing word-only and subword CNNs tests whether partial word spelling helps soften that drop.

\textit{Subword CNN} enriches each word vector with overlapping three-to-five-character fragments inspired by fastText \citep{bojanowski2017}. Each word is wrapped with start and end markers, split into overlapping character chunks, hashed into 8,192 shared buckets, and averaged back into the word vector. A middle-letter swap alters chunks covering the swapped letters, while outer prefix and suffix chunks still activate trained buckets. Both CNNs train with dropout regularization, the AdamW optimizer \citep{adamw}, and validation checkpoint selection.

\subsection{Example of one spelling swap}
In official test article 47, a \emph{Sci/Tech} story about Dell in China, the 5\% noise draw turns \texttt{china consumer pc} into \texttt{china cosnumer pc}. Because \texttt{consumer} appears cleanly a second time later in the story, \textit{Word TF-IDF LR} keeps the word active in its checklist at reduced weight, while neither bigram was in the training dictionary. Table~\ref{tab:observed-swap} shows how all three clean-trained models react to that single typo.

\begin{table}[H]
\centering\normalsize
\begin{tabularx}{\linewidth}{@{}>{\raggedright\arraybackslash}p{32mm}>{\raggedright\arraybackslash}X>{\raggedright\arraybackslash}X@{}}
\toprule
Model & Clean representation (\texttt{consumer}) & Corrupted representation (\texttt{cosnumer}) \\
\midrule
\textit{Word TF-IDF LR} & Two occurrences of \texttt{consumer} activate its word weight, with neither adjacent pair in the vocabulary. & \texttt{cosnumer} is ignored, while the second intact \texttt{consumer} keeps the word active at roughly 40\% lower weight. \\
\addlinespace
\textit{Word CNN} & Maps the position to the learned 96-dimensional vector for \texttt{consumer}. & Replaces the position with the shared unknown-word vector. \\
\addlinespace
\textit{Subword CNN} & Sums the \texttt{consumer} vector with the mean of 16 character-fragment vectors. & Sums the unknown-word vector with the new 16-fragment mean, retaining six shared buckets (including \texttt{<co} and \texttt{mer>}). \\
\bottomrule
\end{tabularx}

\caption[Observed swap representations]{How the clean-trained models represent the \texttt{consumer} $\to$ \texttt{cosnumer} swap in official test article 47. Six of \textit{Subword CNN}'s 16 character fragments keep their pre-swap buckets.}
\label{tab:observed-swap}
\end{table}

\section{Spelling errors and robustness}
Reading studies by \citet{rayner2006} show that transposed middle letters slow human reading somewhat, yet humans still recover meaning easily. Machine translation and classification pipelines enjoy no such margin: \citet{belinkov2018} found heavy quality losses from both natural typos and synthetic swaps even in character-aware neural systems, and \citet{pruthi2019} showed that single-character edits reliably fool classifiers unless paired with explicit word recovery.

One direct defense normalizes spelling before classification, for example by sorting each word's middle letters alphabetically so \texttt{consumer} and \texttt{cosnumer} become identical. Although that trick cancels interior swaps completely, it also merges real English words like \texttt{form} and \texttt{from}, as examined quantitatively in Chapter~\ref{ch:experiments}. Another defense trains on noisy text. Unlike augmentation methods such as EDA \citep{wei2019} that add extra synthetic rows to enlarge a corpus, our noisy arm replaces half of each topic's stories in place so clean and noisy models train on the exact same 20,000-article budget.

Finally, aggregate accuracy and topic-balanced Macro-F1 (defined in Chapter~\ref{ch:method}) give the overall picture, while pairing predictions article by article from clean to noisy text exposes hidden churn: cases where noise breaks a previously correct prediction versus accidental fixes where corrupting a distracting phrase flips a mistake back to the true topic.
